% !TEX root = main.tex
%

%% パッケージ
%ページの余白
\usepackage[a4paper, hmargin=20mm, vmargin=20mm]{geometry}
%XeLaTeXにおける日本語の設定
\usepackage{zxjatype}
\setjamainfont{Noto Serif CJK JP}
\setjasansfont{Noto Sans CJK JP}
\setjamonofont{Noto Sans Mono CJK JP}
%フォント
\usepackage[nomath]{lmodern}
\usepackage{newtxmath}
%図
\usepackage[hiresbb]{graphicx}
\usepackage{caption}
\usepackage[subrefformat=parens]{subcaption}
%数式
\usepackage{amsmath}
\allowdisplaybreaks
\def\equation{\gather}%%hyperref対応
\def\endequation{\endgather}%%hyperref対応
\let\openbox\undefined %newtxmathとamsthmの競合を解消
\usepackage{amsthm}
\newcommand{\hmmax}{0}
\newcommand{\bmmax}{0}
\usepackage{bm}
%アルゴリズム
\usepackage[chapter]{algorithm}
\usepackage{algorithmicx}
\usepackage{algpseudocode}
%URL
\usepackage{url}
\urlstyle{rm}
%引用
\usepackage{cite}%下の3行がない限りhyperrefとは同時に使わない．
\makeatletter
\def\NAT@parse{\typeout{This is a fake Natbib command to fool Hyperref.}}
\makeatother
%索引
\usepackage{makeidx}
\makeindex
%ハイパーリンク
\usepackage{hyperref}%なるべく後
\renewcommand\UrlFont{\rmfamily}
%TODO
\usepackage{todonotes}
\setuptodonotes{inline}
%自分のパッケージ（他のパッケージに依存するから最後にする）
\usepackage[book,amsthm,many]{kmath}
\usepackage{kmacro}

%タイトル
\title{\ktitle}
\author{椛島 健太}
% \西暦
%hypersetup用の設定
\hypersetup{%
    bookmarksnumbered=true,%
    pdftitle={\ktitles},%
    pdfauthor={椛島 健太},%
    pdfsubject={},%
    pdfkeywords={},%
    hidelinks,%
    pdfstartview=FitH,%
    pdfremotestartview=FitH%
}

%参考文献
\def\bibname{参考文献}

%数式の番号（章番号を使うようにする．）
\numberwithin{equation}{chapter}

%行間を日本語向けに変更
\renewcommand{\baselinestretch}{1.3}
